\section{The subtraction}
\label{sec:deamsub}

We can price the stopping right at any volatility we choose.  The
market, meanwhile, hands us one American mid per quote and tells us
nothing else.  This section removes the right from the quote, and the
whole method fits in one line:
\begin{equation}
\boxed{\;\text{find }\sigma^*\ \text{such that}\
A(\sigma^*)=\mathsf{C}(K)\ \text{(or } \mathsf{P}(K)\text{)},\
\text{then report }\sigma^*.\;}
\label{eq:deamroot}
\end{equation}
In words: turn the tree's volatility dial until its \emph{American}
price reproduces the quoted mid, and hand the dial's position to
Part~I as if it had come from a European quote.  This is the standard
desk repair, called \emph{de-Americanization}, and its properties ---
pragmatic, fast, exact inside the model and only there --- were
studied critically by \citet{Burkovska2018}.  The rest of the section
explains why the line works, when it fails, and what each error costs.

\subsection{Why the root is the right volatility}

At first sight \eqref{eq:deamroot} looks like it answers the wrong
question: we wanted the premium subtracted, and instead we matched an
American price.  The resolution is an identity.  At the root, add and
subtract the tree's European leg:
\begin{equation}
\underbrace{\mathsf{C}(K)}_{\text{the quote}}
-\underbrace{\big[A(\sigma^*)-E(\sigma^*)\big]}_{\text{the model's
premium at }\sigma^*}
=E(\sigma^*),
\label{eq:deamcv}
\end{equation}
which holds \emph{exactly}, because $A(\sigma^*)$ equals the quote by
construction.  Read it right to left: the volatility $\sigma^*$ that
reprices the American quote is, at the same time, the volatility at
which the \emph{European} tree price equals the quote minus the
model's premium.  Solving \eqref{eq:deamroot} therefore \emph{is}
subtracting a premium --- the model's --- and the subtraction is
exact to the extent the model's premium matches the market's.  Inside
the model it matches perfectly; outside, the difference between the
true premium and the modelled one (a different diffusion, a different
dividend path) passes straight into $\sigma^*$.  That is the method's
one irreducible approximation, and the chapter's contract
(\cref{sec:deammatters}) states it rather than hiding it.

There is a second, quieter cancellation, and it is why a $1/N_t$
scheme can serve error budgets below a volatility basis point --- the
standard Part~I's fits are held to.  The tree is wrong about
\emph{both} of its legs by discretization; but the two legs share one
lattice, so most of that error is common --- \cref{fig:deamtree}(b)
measured the difference $A-E$ converging
\MacDeamConvRatio$\times$ cleaner than either leg in the median.
This pairing has a name worth knowing: the tree acts as a
\emph{control variate} --- the statistician's term for computing an
error-prone quantity twice, once on each side of a subtraction, so
that the shared error cancels in the difference.  What the pairing
cannot cancel is the European leg's own residual gap against the
analytic Black price, and that gap is what survives into the reported
volatility.  The budget below prices it.

\subsection{When the root exists, and when it does not}

Rooting a scalar equation is legitimate when the function is monotone,
and here it is a theorem: in the Black--Scholes model, the value of an
American option with convex payoff is nondecreasing in the volatility
--- convexity of the payoff propagates to the value function of the
stopping problem, and more volatility then means more value
\citep{ElKarouiJeanblancShreve1998}.  Away from the intrinsic floor
the monotonicity is strict, so when a root of \eqref{eq:deamroot}
exists it is unique.  The theorem is about the continuous problem, not
about any finite tree; the working method leans on something weaker
and verifiable --- \emph{bisection}, which needs only a bracket
$[\sigma_{\text{lo}},\sigma_{\text{hi}}]$ with the quote between
$A(\sigma_{\text{lo}})$ and $A(\sigma_{\text{hi}})$, and halves it
until the root is pinned (the bracket floor sits above
\cref{sec:deamtree}'s drift condition; all constants in
appendix~\ref{app:deamprotocol}).

``When a root exists'' is not decoration.  \Cref{fig:deamroot} draws
the map $\sigma\mapsto A(\sigma)$ for three American puts under a
$4\%$ rate at half a year.  For the out-of-the-money strike ($K=80$)
and the moderately in-the-money one ($K=110$), the curve rises
strictly and crosses its quote --- the dots, priced at a true $25\%$
--- exactly once: the roots are the open circles, and both recover
$25\%$.  The deep strike ($K=160$) is another animal.  Its curve
\emph{starts} on the intrinsic floor of \$\MacDeamRootDeepIntr\ ---
at the true volatility, the converged price carries
\MacDeamRootPlateauGapCents\ cents of time value; the quote \emph{is}
the floor, to the cent --- and only lifts off the floor at
volatilities roughly double the market's.  A quote sitting \emph{at}
that floor is never strictly crossed from above: there is no root,
because there is no volatility in the price to find.  This is
\cref{sec:deamwhere}'s dead zone met in the flesh, and the
output is a flag --- ``this quote determines no volatility'' ---
not a number.  The flag costs nothing downstream: a quote with no
volatility content had nothing to contribute to a volatility fit.

\begin{figure}[htbp]
\centering
\paperfig{\textwidth}{fig_deam_root.pdf}
\caption{The map $\sigma\mapsto A(\sigma)$ for three American puts
($S=100$, $r=4\%$, $t=0.5$).  The $K=80$ and $K=110$ curves cross
their quotes (dots, priced at a true $25\%$) exactly once --- the
open-circled roots.  The $K=160$ curve starts \emph{on} the intrinsic
floor (dashed): at the true volatility the price carries
\MacDeamRootPlateauGapCents\ cents of time value, the quote equals
the floor, and no root exists --- the quote determines no
volatility, and the method says so rather than inventing one.}
\label{fig:deamroot}
\end{figure}

\subsection{From the root to the prepared quote}

One root per strike is also one \emph{premium estimate} per strike:
by \eqref{eq:deamcv}, the model premium at the root is
\begin{equation}
\max\big(\mathsf{C}(K)-E(\sigma^*),\,0\big),
\label{eq:deamhat}
\end{equation}
where the clamp at zero is \cref{prop:deamdom} enforced against
noise --- the true premium cannot be negative, so a negative estimate
is quote or tree noise, not information.  The same single estimate is
subtracted from the bid, the mid, and the ask alike, so the quoted
dollar spread --- the market's own statement of uncertainty, which
Chapter~6 was careful to preserve --- survives the surgery untouched.
The shifted band then flows on exactly as Part~I assumed:
parity conversion to a normalized call (Chapter~2's price in units of
$DF$), inversion to total variance, fitting.  A quote whose root find failed keeps its raw band and an
estimate of zero; the ordinary European sanity filters decide its
fate.

One remark about clocks, before the clocks chapter makes it a theme.
The root $\sigma^*$ is a \emph{calendar-clock} Black volatility: the
tree marched on $t$, where discounting and carry live.  What Part~I
actually consumes is the total variance $w$ --- Chapter~2's
$\sigma^2\tau$ --- inverted from the de-Americanized \emph{price},
and a price is clock-free: only the
reported quotient $\sigma=\sqrt{w/\tau}$ depends on which time
$\tau$ one divides by.  In this book's frozen data the two clocks
agree to the day, so nothing subtle hides in the figures; Chapter~8
is about the cases where they do not.

\subsection{One strike, by the numbers}

Return to the worked board of \cref{fig:deamwedge} and pull out the
$K=90$ call ($S=100$, $r=2\%$, $q_d=6\%$, true $\sigma=25\%$, half a
year; single-quote tree depth in appendix~\ref{app:deamprotocol}).
The tree prices
\[
A=\MacDeamWexAmDollars,
\qquad
E=\MacDeamWexEuDollars,
\qquad
A-E=\MacDeamWexPremiumDollars :
\]
with a $6\%$ yield to capture against a $2\%$ rate, the stopping
right is genuinely valuable on an in-the-money call ---
\MacDeamWexPremiumPct\% of this option's worth sits in the premium.
Hand $A$ to the analytic European inversion and it returns
$\sigma_{\text{naive}}=\MacDeamWexNaivePct\%$ --- a
$+\MacDeamWexBiasBp$ vol bp phantom, the premium booked as
volatility.  Solve \eqref{eq:deamroot} instead and the root comes
back $\sigma^*=\MacDeamWexRootPct\%$: the truth, to bisection
tolerance.

\subsection{The error budget}

The reported $\sigma^*$ carries four error sources; they scale
differently and are worth separating once.
\begin{itemize}[leftmargin=1.6em,itemsep=2pt]
\item \emph{Bisection tolerance.}  The working inversion runs $48$
halvings of a bracket at most $4.0$ wide:
$4\cdot2^{-48}\approx1.4\times10^{-14}$ of volatility.  Nothing.
\item \emph{The finite tree.}  The control-variate pairing protects
the root, but the European-leg gap against analytic Black survives.  Its
measured size on real data: halving the whole-chain depth from
$N_t=256$ to $128$ moves the de-Americanized volatilities of the
frozen running node's fitted quotes by \MacDeamTblDepthMedianBp\ vol
bp in the median and \MacDeamTblDepthMaxBp\ at worst --- so depth is
a numerical-target dial with measurable stakes, not a free speed
knob.  Its small systematic footprint reappears, labelled, in
\cref{sec:deammatters}.
\item \emph{The carry.}  The tree runs on Chapter~6's $(F,D)$ and
dividend model, and $\sigma^*$ inherits them at full strength: the
method removes the early-exercise bias and cannot repair a wrong
forward.  (Circularly, the forward was fitted from American quotes
--- \cref{sec:deamloop} untangles that loop.)
\item \emph{The model.}  The lattice diffusion between dividend
dates, the deterministic dividend path of \cref{sec:deamdivs}, and
the one-sided clamp \eqref{eq:deamhat} are modelling choices.  Exact
in the model, approximate outside it; stated in the contract.
\end{itemize}

We now have the subtraction end to end: a root, its meaning as a
premium removal, its unique-existence theorem, its no-root flag,
and its priced errors.  What remains is the tree's
\emph{carry input} when real cash dividends are on the calendar.
